\section{从 DQN 到 Actor-Critic}

\subsection{Q-Learning 的局限}

\begin{warningbox}{Q-Learning 在连续动作空间的困难}
Q-Learning 需要计算 $\max_{a'} Q(s', a')$，在离散动作空间中可以遍历所有动作。但在连续动作空间中，这个最大化操作变成了一个连续优化问题，计算代价很高。
\end{warningbox}

\subsection{Actor-Critic 方法}

解决方案：同时学习策略（actor）和值函数（critic）：
\begin{itemize}
    \item \textbf{Actor}：$\pi_\theta(a|s)$，参数化策略，输出动作
    \item \textbf{Critic}：$Q_\phi(s, a)$，评估 actor 选择的动作好坏
    \item Actor 用策略梯度更新，critic 用 TD 学习更新
\end{itemize}

\subsection{本章小结}

Actor-Critic 方法结合了策略梯度和值函数方法的优点，是处理连续动作空间的主流范式。

